% Folder 70, batch 2 (archivists' pages 21-40).
% Pass: Opus 5.5 (claude-opus-5-5), 2026-09-29 — first pass, unchecked against the pages by a human.
%
% Pages 26, 27 and 37 carry no mathematics and are absent. Pages 28 and 38 are
% covers in his hand: no \page{}, their words become section headings.
\documentclass[11pt,a4paper]{article}
\input{../preamble/grothendieck}

\folder{70}
\batch{2}
\pages{21}{40}
\foldertitle{Réalisations géométriques de structures combinatoires (n-polyèdres, n-hyperpolyèdres [2-polyèdres réguliers]…) (Vieilles rédactions) : notes manuscrites (s.d.).}
\dating{[à partir de 1976]}
\watermark{Édition de démonstration}

\begin{document}

\page{21}
\note{La page ouvre sur un argument commencé avant ce lot : $G$, $\Gamma_n$, $\Pi$ et $E$ y sont supposés déjà introduits.}
Soit $G$ l'image de $\Gamma_n$ dans $\mathrm{Aut}_{\mathrm{aff}}(E)$ (on suppose maintenant que $k$ est un corps), $\Pi$ le \struck{polyèdre} \add{$n$-\uncertain{carte}} régulière \struck{(\uncertain{partielle})} correspondante. On a donc une représentation \emph{régulière} de $\Pi$. Correspondant : une représentation \emph{fidèle} de $G = \uncertain{\mathrm{Aut}}(\Pi)$ dans $\mathrm{Aut}_{\mathrm{aff}}(E)$, \struck{\ill{}} une application \emph{injective} de $\mathrm{Rep}(\Pi)$ dans $\mathrm{Drap}_{\uncertain{\max}}(E)$. Les applications $D_i(\Pi) \to \struck{\ill{}}\ \mathrm{Grass}_{\mathrm{aff}_i}(E)$ sont-elles aussi injectives -- en d'autres termes, si \struck{$\sigma, \sigma'$ \uncertain{dans}} $g, g' \in G \subset \mathrm{Aut}_{\mathrm{aff}}(E)$ \struck{(\uncertain{engendré} par $\sigma_0,\dots,\sigma_n$) \ill{} \ill{} que $\ill$ \uncertain{produit} qu'il \uncertain{coïncide} \ill{} dimension $i$ $g(f_0, f_1, \dots, f_i)$} \add{\ill{}} sont tels que $g f_i = g' f_i$, $g^{-1}g'$ est-il dans le groupe \struck{engendré} par $\sigma_0, \dots, \sigma_{i-1}, \widehat{\sigma_i}, \sigma_{i+1}, \dots, \sigma_n$ ? Ce n'est pas \uncertain{évident} du tout -- et c'est quelque chose qui n'est pas clair que des résultats sur les \struck{\ill{}} ordres des \uncertain{Longueurs} \uncertain{associés} à la réal. \uncertain{riemannienne} \add{\uncertain{ou non} \struck{des}} $\Pi$, impliquant quelque chose sur les ordres de $\sigma_i\sigma_{i+1}$ !
\note{Deux lignes biffées au milieu de la page, surchargées d'interlignes, ne se lisent que par morceaux ; la question qui suit (« $g^{-1}g'$ est-il dans le groupe \ldots ? ») est nette. « $\mathrm{Grass}_{\mathrm{aff}_i}$ » et « $\mathrm{Drap}_{\max}$ » rendent des abréviations lues « Grass.aff$_i$ » et « Drap.max ».}

\page{22}
\[
\underbrace{\sigma_0\sigma_1}_{\pi}(x_1, x_2, x_3, \dots, x_n) = \bigl(1 + \alpha_0 x_1 - x_2 + \underset{\substack{\| \\ 1+\alpha_1+2\alpha_0+\alpha_0\alpha_1}}{a}\,\Sigma_1,\ x_1 + \underset{\substack{\| \\ 1+\alpha_1}}{b}\,\Sigma_1,\ x_3, \dots\bigr)
\]
\struck{$\pi^{2} = \sigma_1\sigma_0(x_1, x_2, \dots, x_n) = \sigma$.}
\[
\begin{aligned}
\pi^2(x_1, \dots, x_n) = \bigl(&1 + \alpha_0(1 + \alpha_0 x_1 - x_2 + a\Sigma_1) - x_1 - b\Sigma_1 + a\Sigma_1,\\
&1 + \alpha_0 x_1 - x_2 + a\Sigma_1 + b\Sigma_1,\ x_3, \dots\bigr)\\
= \bigl(&(1+\alpha_0) + (\alpha_0^2 - 1)x_1 - \alpha_0 x_2 + (a\alpha_0 + a - b)\Sigma_1,\\
&1 + \alpha_0 x_1 - x_2 + (a+b)\Sigma_1,\ x_3, \dots\bigr)
\end{aligned}
\]
\note{Dans la première ligne de $\pi^2$, les termes sont soulignés deux à deux, pour le regroupement qui suit.}
\[
\begin{aligned}
\pi^3(x_1, \dots, x_n) = \ &1 + \alpha_0(1+\alpha_0) + \alpha_0(\alpha_0^2 - 1)x_1 - \alpha_0^2 x_2 + \alpha_0(a\alpha_0 + a - b)\Sigma_1\\
&-1\ \struck{\ill{}}\quad - \alpha_0 x_1 \quad + x_2 \quad - (a+b)\Sigma_1 + a\Sigma_1\\
&\uncertain{(1+\alpha_0) + (\alpha_0^2-1)x_1 - \alpha_0 x_2 + [a(\alpha_0 + 1)]\Sigma_1,\ x_3 \ldots}\\
= \ &\struck{(1+}\ \alpha_0(1+\alpha_0) + \alpha_0(\alpha_0^2 - 2)x_1 + (1 - \alpha_0^2)x_2 + \struck{(a\alpha_0^2 + a\alpha_0 \ldots}\\
&(\alpha_0+1)(a\alpha_0 - b)\ \struck{\ill{}}\ \Sigma_1,\ \ldots
\end{aligned}
\]
\note{La troisième ligne, d'une encre plus pâle et en partie barrée d'obliques, est reliée par une accolade et une flèche au crayon à la suite ; la fin du calcul est au crayon.}

$\alpha_0 = -1$
\[
\begin{cases}
\alpha_0(1 + \alpha_0) = 0 \qquad \alpha_0(\alpha_0^2 - 2) = \struck{1}\ \uncertain{-1}? \qquad \alpha_0^2 = 1 \quad \big|\ \alpha_0 = -1\\
\struck{\ill{}}
\end{cases}
\]

\page{23}
\struck{$u(x,y) = \sigma_0\sigma_1(x,y) =$ \quad $\sigma_0(x,y) = (1 - (1+\eta_0)x + \alpha_0 y,\ y)$, \quad $\sigma_1(x,y) = (-\eta_1 y,\ x + (\eta_1+1)y)$}

\struck{$u(x,y) = \sigma_0\sigma_1(x,y) =$}
\[
\begin{aligned}
\sigma_0(x_1, \dots, x_{n+1}) &= (1 - x_1 + \alpha_0\Sigma_0,\ x_2, \dots, x_{n+1})\\
\sigma_i(x_1, \dots, x_{n+1}) &= (x_1, \dots, x_{i-1},\ x_{i+1} - (1+\alpha_i)\Sigma_i,\ x_i + (1+\alpha_i)\Sigma_i,\ x_{i+2}, \dots, x_{n+1}) \quad (1\le i\le n-1)\\
\sigma_n(x_1, \dots, x_{n+1}) &= (x_1, \dots, x_{n-1},\ x_{n+1},\ x_n)
\end{aligned}
\]
\[
\struck{f(x_1, \dots, x_{n+1}) = \sum_{1\le i\le n} a_{ii}\, x_i^2 + \sum_{1\le i<j\le n+1} a_{ij}\, x_i x_j + \sum_{1\le i\le n+1} c_i x_i}
\]

\struck{\uncertain{Invariants} ?}
\[
\begin{array}{ll}
H_0 : & 2x_1 \struck{=} - \alpha_0(x_2 + \cdots + x_{n+1}) = z\\
D_0 : & x_2 = x_3 = \cdots = x_{n+1} = z = 0 \qquad (\text{pt à l'infini de la } \struck{\ill{}}\ \mathrm{dr}(s_0 s_1))\\[4pt]
H_i : & x_i - x_{i+1} + (1+\alpha_i)(\struck{x_{i+2} =}\ x_{i+2} + \cdots + x_{n+1}) = 0\\
D_i : & x_1 = \cdots = x_{i-1} = x_{i+2} = \cdots = x_{n+1} = z = 0
\end{array}
\]
\note{Tout le cadre $H_0, D_0, H_i, D_i$ est barré de trois longues obliques ; il est repris proprement p.~24. Le signe entre $2x_1$ et $\alpha_0$ est surchargé.}

\[
\begin{aligned}
\overline{\sigma}_0(\underbrace{x_1, \dots, x_{n+1}, \overset{x_0}{\overset{\|}{z}}}_{\overline{x}}) &= (x_0 - x_1 + \alpha_0(x_2 + \cdots + x_{n+1}),\ x_2, \dots, x_{n+1},\ x_0)\\
&= \overline{x} + \struck{\ill{}}\ [e'_0 - 2e'_1 + \alpha_0(e'_2 + \cdots + e'_{n+1})] \otimes e_1\\
\overrightarrow{\sigma_i}(\overline{x}) &= \overline{x} + [e'_i - e'_{i+1} + (1+\alpha_i)(e'_{i+2} + \cdots + e'_{n+1})] \otimes (-e_i + e_{i+1}) \quad (1\le i\le n)\\
\overrightarrow{\sigma_n}(\overline{x}) &= \overline{x} + (e'_n - e'_{n+1}) \otimes (-e_n + e_{n+1})
\end{aligned}
\]
\note{Dans la ligne de $\overline{\sigma}_0$, le coefficient de $e'_1$ est surchargé : un $1$ récrit, lu $2$ d'après $H_0$ ; le signe $\otimes$ devant $e_1$ est de même empâté. Les barres sur $\sigma_i$, $\sigma_n$ portent une pointe de flèche.}

\page{24}
\[
\begin{array}{lll}
H_0 : & \overset{\substack{z\ \add{\text{variable d'homogénéité}}\\ \|}}{\struck{x_0}} - 2x_1 + \alpha_0(x_2 + \cdots + x_{n+1}) = 0 &\\
\delta_0 : & \text{pt à l'infini dans la direction de } e_1 = s_1 - s_0 &\\[4pt]
H_i : & x_i - x_{i+1} + (1+\alpha_i)(x_{i+2} + \cdots + x_{n+1}) = 0 & (1\le i\le n)\\
\delta_i : & \text{pt à l'infini dans la direction de } e_{i+1} - e_i = s_{i+1} - s_i &
\end{array}
\]
\note{La glose « $z$ variable d'homogénéité », au crayon, est écrite au-dessus de $x_0$ ; le $z$ est surchargé en $2$ ou l'inverse, lu $z$ d'après la p.~23. Sous ces quatre lignes, un trait ; le reste de la page est blanc.}

\page{25}
\[
\sigma_0 s_i = -(1+\alpha_i)s_0 + (1+\alpha_i)s_1 + s_i \qquad i\ge 2
\]
\[
\sigma_0 s_{i+1} = \sigma_0\sigma_i s_i = \sigma_i\sigma_0 s_i = -(1+\alpha_i)s_0 + (1+\alpha_i)s_1 + s_{i+1}
\]
\note{Ainsi sur la page : les indices de $\alpha$ dans ces deux lignes ne suivent pas ceux de $s$.}
\[
\begin{aligned}
\sigma_i s_j &= s_j \quad 0\le j\le i-1\\
\sigma_i s_i &= s_{i+1}\\
\sigma_i s_{i+1} &= s_i\\
\sigma_i s_{i+2} &= (\underbrace{0, \dots, 0}_{i-1},\ -(1+\alpha_i),\ (1+\alpha_i),\ 1,\ 0, 0, \dots) = s_{i+2} + (1+\alpha_i)(s_{i+1} - s_i)\\
&= -(1+\alpha_i)s_i + (1+\alpha_i)s_{i+1} + s_{i+2}
\end{aligned}
\]

\struck{$\sigma_i s_{i+2}$} \quad \struck{\ill{}}\ pour $0\le i\le n$
\[
\boxed{\sigma_i s_j = \begin{cases}
s_j & \text{si } 0\le j\le i-1\\
s_{i+1} & \text{si } j = i\\
s_i & \text{si } j = i+1\\
-(1+\alpha_i)s_i + (1+\alpha_i)s_{i+1} + s_j & \text{si } i+2\le j\le n+1
\end{cases}}
\]
\[
\struck{\sigma_i s_j - s_j =} \quad \sigma_i - \mathrm{id} = \ell_i \otimes (s_{i+1} - s_i), \qquad \ell_i = s'_i - s'_{i+1} + (1+\alpha_i)(s'_{i+2} + \cdots + s'_{n+1})
\]
\struck{$\sigma_i x = x +$}
\[
\sigma_i s_j - s_j = \begin{cases}
0 & \text{si } 0\le j\le i-1\\
s_{i+1} - s_i & \text{si } j = i\\
-(s_{i+1} - s_i) & \text{si } j = i+1\\
(1+\alpha_i)(s_{i+1} - s_i) & \text{si } j\ge i+2
\end{cases}
\]

\section{$n$-quasi-polyèdres (vieilles rédactions)}
\note{Titre de sa main, au crayon, sur la feuille de couverture p.~28 : « $n$-quasi-polyèdres (vieilles rédactions) ».}

\page{29}
\emph{$n$-Quasi-polyèdres réguliers} en \uncertain{coordonnées} \ill{}
\[
(1)\quad \begin{cases}
\sigma_0,\dots,\sigma_n \in \mathrm{Aff}(E), & E \text{ affine de rang } n+1\\
\sigma_j\sigma_i = \sigma_i\sigma_j & \text{si } 0\le i,\ \struck{i+1}\ i+2\le j\le n
\end{cases}
\]
\[
(2)\quad E_0 = \{s_0\} \subset E_1 \subset \cdots \subset E_n \quad \text{drapeau de } E
\]
\[
(3)\quad \begin{cases}
\sigma_i E_j = E_j & \text{si } j\ne i\\
\sigma_i E_i \ne E_i
\end{cases}
\qquad
(4)\quad \begin{cases}
s_0\\
s_1 = \sigma_0 s_0\\
\cdots\\
s_i = \sigma_{i-1} s_{i-1} & (1\le i\le n+1)\\
\cdots\\
s_{n+1} = \sigma_n s_n
\end{cases}
\]

\emph{Prop} $(s_0,\dots,s_{n+1})$ est une base affine de $E$, et
\[
\begin{aligned}
(5)&\quad E_i = \struck{\ill{}}\ \mathrm{Env}(s_0,\dots,s_i)\\
(6)&\quad \sigma_k s_i = s_i \quad \text{si } 0\le i< k\le \struck{n+1}\ n
\end{aligned}
\]
\add{pour \uncertain{tout} $i$, $0\le i\le n+1$,}
\note{L'abréviation que nous rendons par $\mathrm{Env}$ (l'engendré affine, semble-t-il) est lue « Env » à chaque occurrence de la page, sans certitude sur les deux dernières lettres. Dans (1), le « $0\le i$ » est suivi d'un $i+1$ biffé.}

Par induction \add{sur $i$, on prouve que $(s_0,\dots,s_i)$ est affinement libre et} (5)$_i$ (6)$_i$. Trivial si $i=0$. Supposons \add{$i\ge 1$ et} prouvé pour $j<i$. Je \uncertain{dis} \ill{} :
\[
\begin{cases}
s_i \notin E_{i-1} \ (= \mathrm{Env}(s_0,\dots,s_{i-1}) \text{ par hyp.\ réc.})\\
E_i = \mathrm{Env}(E_{i-1}, s_i)\\
\sigma_k s_i = s_i \quad i<k\le n
\end{cases}
\]
Si on avait $s_i \in E_{i-1}$, i.e.\ $\sigma_{i-1}s_{i-1} \in E_{i-1}$, comme $\sigma_{i-1}s_j = s_j$ pour $0\le j\le i-2$ et $E_{i-1} = \mathrm{Env}(s_0,\dots,s_{i-1})$ par réc., on aurait
\[
\sigma_{i-1}E_{i-1} = \mathrm{Env}(\underbrace{\sigma_{i-1}s_0}_{=s_0},\dots,\underbrace{\sigma_{i-1}s_{i-2}}_{=s_{i-2}},\underbrace{\sigma_{i-1}s_{i-1}}_{\in E_{i-1}}) \subset E_{i-1}
\]
donc $\sigma_{i-1}E_{i-1} = E_{i-1}$, contrairement à (3 b), OK.

On a, comme $\sigma_{i-1}E_i = E_i$ et $s_{i-1} \in E_{i-1}$ (hyp.\ réc.) $\subset E_i$, $\sigma_{i-1}(s_{i-1}) \subset E_i$, i.e.\ $s_i \in E_i$. Comme $s_i \notin E_{i-1}$, on a donc $E_i = \mathrm{Env}(E_{i-1}, s_i)$. Enfin, pour \struck{\ill{}} $i<k\le n$,
\[
\sigma_k s_i = \sigma_k\sigma_{i-1}s_{i-1} = \sigma_{i-1}(\sigma_k s_{i-1}) \overset{\text{réc}}{=} \sigma_{i-1}(s_{i-1}) = s_i.
\]
OK.

\page{30}
On utilise les $s_i$ ($0\le i\le n+1$) comme base de $E$, \ill{} \ill{} \uncertain{comme origine}, et les
\[
(7)\quad e_i = s_i - s_0 \qquad 1\le i\le n+1
\]
comme base de l'espace des translations $V$ de $E$. On va expliciter les matrices (à $n+2$ lignes et $n+2$ colonnes) des $\sigma_i$, $0\le i\le n$. Revient à la connaissance des $\sigma_i s_j$, $0\le i\le n$, $0\le j\le n+1$. Je dis qu'il suffit de connaître
\[
(8)\quad \underset{\textstyle \sigma_i^2 s_i}{\sigma_i s_{i+1}} = \struck{u_{i+1}}\ u_i, \qquad \sigma_i s_{i+2} = v_i
\]
Car on a
\[
(9)\quad \begin{cases}
\sigma_i s_j = s_j & \text{si } 0\le j< i\\
\sigma_i s_i = s_{i+1}\\
\sigma_i s_{i+1} = u_i \in E_{i+1} & E_{i+1} = \mathrm{Env}(E_{i-1}, s_{i+1}, u_i)\\
\sigma_i s_{i+2} = v_i \in E_{i+2} & E_{i+2} = \mathrm{Env}(E_{i+1}, v_i) \quad (0\le i\le n-1)\\
\sigma_i s_j = \sigma_{j-1}\cdots\sigma_{i+2}(v_i) & \text{si } \uncertain{i+2< j\le n+1}
\end{cases}
\]
\note{Une flèche renvoie la condition « $(0\le i\le n-1)$ » à la ligne de $v_i$. La condition sur $j$ dans la dernière ligne est surchargée.}

Car $s_j = \sigma_{j-1}\cdots\sigma_{i+2}\,s_{i+2}$, donc
\[
\sigma_i s_j = \sigma_i\sigma_{j-1}\cdots\sigma_{i+2}\,s_{i+2} = \sigma_{j-1}\cdots\sigma_{i+2}\,\underbrace{\sigma_i s_{i+2}}_{v_i}
\]
Posons \struck{\ill{} $= s_0+$}
\[
(\ill)\quad \begin{cases}
u_i = (u_{i,1},\dots,u_{i,i+1}) = s_0 + \sum_{1\le k\le i+1} u_{i,k}\,e_k\\
v_i = (v_{i,1},\dots,v_{i,i+2}) = s_0 + \sum_{1\le k\le i+2} v_{i,k}\,e_k
\end{cases}
\]
\ill{} \ill{} \uncertain{obtient} si $j\ge i+3$
\[
\begin{aligned}
\sigma_i s_j &= \sigma_{j-1}\cdots\sigma_{i+2}\,(s_0 + v_{i,1}e_1 + \cdots + v_{i,i+1}e_{i+1} + v_{i,i+2}e_{i+2})\\
&= s_0 + \sum_{1\le k\le i+1} v_{i,k}\,e_k + v_{i,i+2}\,e_j
\end{aligned}
\]
i.e.
\[
(10)\quad \begin{cases}
\sigma_i s_j = (v_{i,1},\dots,v_{i,i+1},\ 0,\dots,0,\ \overset{j}{v_{i,i+2}},\ 0,\dots,0)\\
\qquad = s_0 + \sum_{1\le k\le i+1} v_{i,k}\,e_k + v_{i,i+2}\,e_j
\end{cases}
\quad \text{si } \struck{\ill{}}\ i+2\le j\le n+1
\]

\page{31}
compte tenu que, pour \uncertain{$0\le i\le n-1$}, $\sigma_i e_i = e_{i+1}$. On a \add{plus \uncertain{explicitement}}
\[
(11)\quad
\begin{cases}
\sigma_i e_j = e_j & \text{si } 1\le j< i\\
\sigma_i e_i = e_{i+1} \quad \struck{(\text{si } i\ge 1)} & \\
\sigma_i e_{i+1} = \struck{\ldots}
  \begin{cases} u_0 - s_1 = (u_0-s_0) - e_1 = (u_{01}-1)\,e_1 & \text{si } i=0\\
  u_i - s_0 = \sum_{1\le k\le i+1} u_{ik}\, e_k & \text{si } i\ge 1\end{cases}\\
\struck{\sigma_i e_{i+2} = \ldots}\\
\Bigl[\sigma_0 e_{i+2} = \Bigr]
  \begin{cases} v_0 - s_1 = (v_0-s_0) - e_1 = (v_{01}-1)\,e_1 + v_{0,2}\, e_2 & \text{si } i=0\\
  v_i - s_0 = \sum_{1\le k\le i+2} v_{i,k}\, e_k & \text{si } i\ge 1\end{cases}\\
\sigma_i e_j = \begin{cases} (v_{01}-1)\,e_1 + v_{0j}\, e_j & \text{si } i=0\\
  \sum_{1\le k\le i+1} v_{ik}\, e_k + v_{i,i+2}\, e_j & \text{si } i\ge 1\end{cases}
  & \text{si } i+2\le j\le n+1
\end{cases}
\]
\note{« plus \uncertain{explicitement} » est écrit sous « compte tenu », en interligne. La ligne $\sigma_i e_{i+2}$ est biffée et récrite, entre crochets, en $\sigma_0 e_{i+2}$ ; les bornes des sommes sont lues sur des indices très petits.}

\struck{Conditions}

Points fixes de $\sigma_0$

\struck{$\sigma_0 =$}
\[
(12)\quad \sigma_0 =
\begin{pmatrix}
u_{01}-1 & v_{01}-1 & v_{01}-1 & v_{01}-1 & \cdots & v_{01}-1 & 1\\
0 & v_{02} & 0 & 0 & \cdots & 0 & 0\\
0 & 0 & v_{02} & 0 & \cdots & 0 & 0\\
0 & 0 & 0 & v_{02} & \cdots & 0 & 0\\
\vdots & & & & \ddots & & \vdots\\
0 & 0 & 0 & 0 & \cdots & v_{02} & 0\\
0 & 0 & 0 & 0 & \cdots & 0 & 1
\end{pmatrix}
\]
\note{Les colonnes sont numérotées $1, 2, \dots, n+1, n+2$ en tête, les lignes $1, 2, \dots, n+1, n+2$ en marge gauche ; la dernière colonne est celle de la translation.}

\[
(13)\quad \sigma_0(x_1,\dots,x_{n+1}) = \bigl(1 + (u_{01}-1)\,x_1 + (v_{01}-1)\,\Sigma_0,\ v_{02}\,x_2,\ \dots,\ v_{02}\,x_{n+1}\bigr),
\qquad \Sigma_0 = x_2 + \cdots + x_{n+1}
\]
\[
(14)\quad H_0 : \begin{cases}
1 + (u_{01}-2)\,x_1 + (v_{01}-1)\,\Sigma_0 = 0\\
(v_{02}-1)\,x_2 = \cdots = (v_{02}-1)\,x_{n+1} = 0
\end{cases}
\]
C'est une réflexion projective sss les \uncertain{dernières} relations sont conséquences de la première \uncertain{homogénéisée}, ou encore sss $\boxed{v_{02} = 1}$. Je vais \uncertain{prendre} \ill{} $\sigma_0$ les paramètres \uncertain{libres} $u_{01}, v_{01}$, \ill{} \uncertain{$\lambda_0$ et $\mu_0$} :
\[
(15)\quad \sigma_0(x_1,\dots,x_{n+1}) = \bigl(1 + (\lambda_0-1)\,x_1 + (\mu_0-1)\,\Sigma_0,\ x_2,\ \dots,\ x_{n+1}\bigr),
\qquad \Sigma_0 = x_2+\cdots+x_{n+1}
\]

\page{32}
\[
(16)\quad \sigma_i =
\left(\begin{array}{ccc|c|ccccc|c}
1 & & 0 & 0 & u_{i1} & v_{i1} & v_{i1} & \cdots & v_{i1} & 0\\
 & \ddots & & \vdots & \vdots & \vdots & \vdots & & \vdots & \vdots\\
0 & & 1 & 0 & u_{i,i-1} & v_{i,i-1} & v_{i,i-1} & \cdots & v_{i,i-1} & 0\\
 & & & 0 & u_{i,i} & v_{i,i} & v_{i,i} & \cdots & v_{i,i} & 0\\
 & 0 & & 1 & u_{i,i+1} & v_{i,i+1} & v_{i,i+1} & \cdots & v_{i,i+1} & 0\\
 & & & 0 & 0 & v_{i,i+2} & 0 & \cdots & 0 & 0\\
 & & & 0 & 0 & 0 & v_{i,i+2} & & 0 & 0\\
 & & & \vdots & \vdots & & & \ddots & & \vdots\\
 & & & 0 & 0 & 0 & 0 & & v_{i,i+2} & 0\\
\hline
0 & \cdots & 0 & 0 & 0 & 0 & 0 & \cdots & 0 & 1
\end{array}\right)
\]
\note{Matrice dessinée par blocs : un bloc identité de taille $i-1$ en haut à gauche, des blocs nuls figurés par de grands zéros ovales ; la disposition ci-dessus suit le dessin sans le reproduire exactement. Plusieurs indices du bloc central sont repassés et incertains.}

\[
(17)\quad \begin{aligned}
\sigma_i(x_1,\dots,x_{n+1}) = \bigl(&x_1 + u_{i1}\,x_{i+1} + v_{i1}\,\Sigma_i,\ x_2 + u_{i2}\,x_{i+1} + v_{i2}\,\Sigma_i,\ \dots\\
&\dots,\ x_{i-1} + u_{i,i-1}\,x_{i+1} + v_{i,i-1}\,\Sigma_i ;\\
&u_{i,i}\,x_{i+1} + v_{i,i}\,\Sigma_i,\ x_i + u_{i,i+1}\,x_{i+1} + v_{i,i+1}\,\Sigma_i ;\\
&v_{i,i+2}\,x_{i+2},\ \dots,\ v_{i,i+2}\,x_{n+1}\bigr),
\end{aligned}
\qquad \Sigma_i = x_{i+2} + \cdots + x_{n+1}
\]
\[
(18)\quad H_i = \begin{cases}
u_{i1}\,x_{i+1} + v_{i1}\,\Sigma_i = \uncertain{u_{i2}\,x_{i+1} + v_{i2}\,\Sigma_i} = \cdots = u_{i,i-1}\,x_{i+1} + v_{i,i-1}\,\Sigma_i = 0\\
\add{-x_i+}\ u_{ii}\,x_{i+1} + v_{ii}\,\Sigma_i = 0,\quad x_i + (u_{i,i+1}-1)\,x_{i+1} + v_{i,i+1}\,\Sigma_i = 0\\
(v_{i,i+2}-1)\,x_{i+2} = \cdots = (v_{i,i+2}-1)\,x_{n+1} = 0
\end{cases}
\]

$\sigma_i$ est une réflexion sss ces relations sont conséquences de la $i$-ième (qui exprime $x_i$ \uncertain{comme combinaison linéaire} de $x_{i+1},\dots,x_{n+1}$). Ceci signifie :
\[
(19)\quad \begin{cases}
u_{ii} + u_{i,i+1} = 1, & v_{i,i} + v_{i,i+1} = 0\\
v_{i,i+2} = 1\\
u_{i1} = \cdots = u_{i,i-1} = 0, & v_{i1} = \cdots = v_{i,i-1} = 0
\end{cases}
\]
donc les seuls paramètres qui restent sont $u_{i,i} = \lambda_i$ et $v_{i,i} = \mu_i$ \struck{(\uncertain{définis pour} $0\le i\le n-1$, pas pour $\sigma_n$)}, (\ill{})
\[
(20)\quad \sigma_i(x_1,\dots,x_{n+1}) = \bigl(x_1,\dots,x_{i-1},\ \lambda_i\,x_{i+1} + \mu_i\,\Sigma_i,\ x_i + (1-\lambda_i)\,x_{i+1} - \mu_i\,\Sigma_i,\ x_{i+2},\dots,x_{n+1}\bigr),
\qquad \Sigma_i = x_{i+2}+\cdots+x_{n+1}
\]

\page{33}
On a ($1\le i\le n$)
\[
\sigma_i^2(x_1,\dots,x_{n+1}) = (x_1,\dots,x_{i-1};\ X_i,\ X_{i+1},\ x_{i+2},\dots,x_{n+1})
\]
\[
\begin{aligned}
X_i &= \lambda_i\bigl[x_i + (1-\lambda_i)\,x_{i+1} - \mu_i\,\Sigma_i\bigr] + \mu_i\,\Sigma_i = \lambda_i\,x_i + \lambda_i(1-\lambda_i)\,x_{i+1}\\
X_{i+1} &= \bigl[\lambda_i\,x_{i+1} + \mu_i\,\Sigma_i\bigr] + (1-\lambda_i)\bigl[x_i + (1-\lambda_i)\,x_{i+1} - \mu_i\,\Sigma_i\bigr] - \mu_i\,\Sigma_i\\
&= (1-\lambda_i)\,x_i + x_{i+1} + (1-\lambda_i)\,\mu_i\,\Sigma_i
\end{aligned}
\]
\note{Des traits obliques relient les termes en $\mu_i\Sigma_i$ qui se compensent. Le coefficient de $x_{i+1}$ dans $X_{i+1}$ est écrit $1$, comme sur la page.}

On a $X_i \equiv x_i$ sss $\lambda_i = 1$, $X_{i+1} \equiv x_{i+1}$ sss $\lambda_i = 1$ \quad
$(21)\ \boxed{\sigma_i^2 = \mathrm{id} \text{ sss } \lambda_i = 1}$

\ill{}
\[
\sigma_0^2(x_1,\dots,x_{n+1}) = (X_1,\ x_2,\ \dots,\ x_{n+1})
\]
\[
\begin{aligned}
X_1 &= 1 + (\lambda_0-1)\bigl[1 + (\lambda_0-1)\,x_1 + (\mu_0-1)\,\Sigma_0\bigr] + (\mu_0-1)\,\Sigma_0\\
&= \lambda_0 + (\lambda_0-1)^2\,x_1 + \lambda_0(\mu_0-1)\,\Sigma_0
\end{aligned}
\]
$X_1 \equiv x_1$ sss $\lambda_0 = 0$, i.e. \quad $(22)\ \boxed{\sigma_0^2 = \mathrm{id} \text{ sss } \lambda_0 = 0}$

\begin{itemize}
\item \emph{Cas général} : paramètres $\lambda_i$ ($0\le i\le n$), $\mu_i$ ($0\le i\le n-1$), donc \emph{$2n+1$ paramètres}.
\item \emph{Cas $n$-polyèdral} ($\sigma_i^2 = 1$ pour $0\le i\le n$) : $\lambda_0 = 0$, $\lambda_i = 1$ ($1\le i\le n$) ; il reste $n$ paramètres $\mu_i$\,\ldots
\end{itemize}

\page{34}
\[
\begin{cases}
\sigma_0(x,y) = \bigl(1 + (\lambda_0-1)\,x + (\mu_0-1)\,y,\ y\bigr) = \bigl(1 + (\lambda_0-1)\,x + \alpha_0\,y,\ y\bigr) = (1 + \eta_0\,x + \alpha_0\,y,\ y)\\
\sigma_1(x,y) = \bigl(\lambda_1\,y,\ x + (1-\lambda_1)\,y\bigr)
\end{cases}
\]
$k$ corps de car.\ $p \ge 0$.

Partie homogène de $\sigma_0$ : $\begin{pmatrix}\eta_0 & \alpha_0\\ 0 & 1\end{pmatrix}$

\struck{$\det(t\,\mathrm{id} - \sigma_{0v}) =$} valeurs propres : $1$ et $\lambda_0 - 1 = \eta_0$.

\struck{$\sigma_0$}
\[
\begin{pmatrix}a & c\\ 0 & b\end{pmatrix}\begin{pmatrix}a' & c'\\ 0 & b'\end{pmatrix}
= \begin{pmatrix}aa' & ac' + b'c\\ 0 & bb'\end{pmatrix}
\overset{\text{si } b'=1}{=} \begin{pmatrix}aa' & c + ac'\\ 0 & b\end{pmatrix}
\]
\[
\begin{pmatrix}a & c\\ 0 & b\end{pmatrix}\begin{pmatrix}\eta_0 & \alpha_0\\ 0 & 1\end{pmatrix}
= \begin{pmatrix}\eta_0\,a & c + \alpha_0\,a\\ 0 & b\end{pmatrix}
\]
\[
\begin{pmatrix}\eta_0 & \alpha_0\\ 0 & 1\end{pmatrix}^2 = \begin{pmatrix}\eta_0^2 & \alpha_0 + \uncertain{\alpha_0\eta_0}\\ 0 & 1\end{pmatrix}
\qquad
\begin{pmatrix}\eta_0 & \alpha_0\\ 0 & 1\end{pmatrix}^3 = \begin{pmatrix}\eta_0^3 & \alpha_0(1 + \eta_0 + \eta_0^2)\\ 0 & 1\end{pmatrix}
\]
\[
\begin{pmatrix}\eta_0 & \alpha_0\\ 0 & 1\end{pmatrix}^n = \begin{pmatrix}\eta_0^n & \alpha_0(1 + \eta_0 + \cdots + \eta_0^{n-1})\\ 0 & 1\end{pmatrix}
\]
\note{Le second coefficient du carré est surchargé et en partie biffé.}

\[
\sigma_{0v}^n = 1 \iff
\begin{cases}
\eta_0^n = 1 & \text{i.e. } (\eta_0 - 1)(1 + \eta_0 + \cdots + \eta_0^{n-1}) = 0\\
\struck{\alpha_0 \ldots} & \alpha_0(1 + \eta_0 + \cdots + \eta_0^{n-1}) = 0
\end{cases}
\]
$\Updownarrow$

$\sigma_0^n$ une translation. Deux cas :
\begin{itemize}
\item[a)] $1 + \eta_0 + \cdots + \eta_0^{n-1} = 0$, i.e. \struck{\ill{}} $\eta_0^n = 1$ et ($\eta_0 \ne 1$ ou $p \nmid n$) ;
\item[b)] $\eta_0 = 1$, $\alpha_0 = 0$ [i.e.\ $\sigma_{0v} = \mathrm{id}$, i.e.\ $\sigma_0$ est la translation par $(1,0)$] et $p \mid n$.
\end{itemize}
\note{Dans a), la parenthèse est très surchargée : « $\eta_0 \ne 1$ ou $p\nmid n$ » est une lecture, le signe de non-divisibilité étant un trait épais.}

$\sigma_0^n(0) = {?}$ \qquad \struck{$\sigma_0(0) = (1,0)$, $\sigma_0^2(0) = (1+\eta_0 \ldots$}
\[
\begin{aligned}
\sigma_0(x,0) &= (1 + \eta_0\,x,\ 0)\\
\sigma_0^2(x,0) &= 1 + \eta_0 + \eta_0^2\,x\\
\sigma_0^3(x,0) &= 1 + \eta_0 + \eta_0^2 + \eta_0^3\,x\\
&\cdots\\
\sigma_0^n(x,0) &= 1 + \eta_0 + \cdots + \eta_0^{n-1} + \eta_0^n\,x =
\begin{cases}
(\text{cas a}) & \eta_0^n\,x = x\\
(\text{cas b}) & \text{translation par } (n,0) \text{ (qui est } \ill \text{ par l'} \ill\text{)}
\end{cases}
\end{aligned}
\]

\page{35}
On trouve donc
\[
\sigma_0^{\varepsilon} = \mathrm{id} \iff 1 + \eta_0 + \cdots + \eta_0^{\varepsilon-1} = 0
\]
donc

ordre de $\sigma_0$ = ordre de $\eta_0 \in k^*$ \emph{sauf} si $\eta_0 = 1$ \add{i.e.\ $\det\sigma_{0v} = 1$}, auquel cas $\sigma_{0v}$ est la transvection $\begin{pmatrix}1 & \alpha_0\\ 0 & 1\end{pmatrix}$ (qui est d'ordre $1$ sss $\alpha_0 = 0$, \uncertain{sinon} d'ordre $p$)
\[
\eta_0 = 1\
\begin{cases}
\text{a)}\ \alpha_0 = \uncertain{1} & \sigma_0 = \mathrm{id}\quad \emph{cas impropre}\\[2pt]
\text{b)}\ \alpha_0 \ne 1 & \begin{cases}
p = 0 & \sigma_{0v} \text{ d'ordre infini, a fortiori } \sigma_0 \text{ d'ordre inf.}\\
p > 0 & \sigma_{0v} \text{ d'ordre } p,\ \sigma_0^p\ \struck{= \text{translation}}\ \ldots
\end{cases}
\end{cases}
\]
\note{Dans a) et b), le chiffre après $\alpha_0$ est repassé en encre plus foncée et se lit $1$ ; la ligne précédente dit « d'ordre $1$ sss $\alpha_0 = 0$ ». Sous la ligne $p>0$, un passage biffé de hachures obliques : \struck{\ill{} $(p,0)$ \ill{} ; $p\mid n \Rightarrow \sigma_0^p = \mathrm{id}$ ; $p\nmid n$ \ldots\ $\sigma_0$ d'ordre $p^2$}.}

\[
\begin{aligned}
\sigma_1^2(x,y) &= \bigl(\lambda_1\,x + \lambda_1(1-\lambda_1)\,y,\ (1-\lambda_1)\,x + (\lambda_1 + (1-\lambda_1)^2)\,y\bigr)\\
&= \bigl(\lambda_1\,x + \lambda_1(1-\lambda_1)\,y,\ (1-\lambda_1)\,x + (\lambda_1^2 - \lambda_1 + 1)\,y\bigr)
\end{aligned}
\]
\[
\sigma_1 = \begin{pmatrix}0 & \lambda_1\\ 1 & 1-\lambda_1\end{pmatrix}
\qquad
\det(t\,\mathrm{id} - \sigma_1) = t^2 - (1-\lambda_1)\,t - \lambda_1 = \struck{t^2 - (1+\delta)\,t + \delta}
\]
Valeurs propres \struck{$\eta, \eta'$} : $1$, $-\lambda_1 = \delta_1$.

\struck{$-\lambda_1 = \eta\eta'$, $1 - \lambda_1 = \eta + \eta'$ ; $-\lambda_1 = \delta_1 = \delta$ ; $\eta\eta' = \delta$, $\eta + \eta' = 1+\delta$ ; donc $\eta^n = \eta'^n = 1 \Rightarrow \delta^n = 1$ ; $\eta/\delta$, $\eta'/\delta$ \ldots\ ; $x = \lambda_1 y$}
\note{Tout ce bloc est biffé de traits obliques et entouré ; les cadres et flèches qui en relient les morceaux ne sont pas reproduits.}

Ordre de $\sigma_1$ = ordre de $\delta_1 \in k^*$, \emph{sauf si $\delta = 1$}, i.e.\ $\lambda_1 = -1$ :
$\sigma_1(x,y) = (-y,\ x + 2y)$ transvection \struck{\ill{}} \uncertain{identité}, d'ordre infini si $p = 0$, d'ordre $p$ sinon.

\page{36}
\[
\begin{cases}
\sigma_0(x_1,\dots,x_{n+1}) = (1 + \eta_0\,x_1 + \alpha_0\,\Sigma_0,\ x_2,\ \dots,\ x_{n+1})\\
\cdots\\
\sigma_i(x_1,\dots,x_{n+1}) = \bigl(x_1,\dots,x_{i-1},\ -\eta_i\,x_{i+1} + (1+\alpha_i)\,\Sigma_i,\ x_i + (\eta_i + 1)\,x_{i+1} - (1+\alpha_i)\,\Sigma_i,\ x_{i+2},\dots,x_{n+1}\bigr) & (1\le i\le n-1)\\
\cdots\\
\sigma_n(x_1,\dots,x_{n+1}) = \bigl(x_1,\ \dots,\ x_{n-1},\ -\eta_n\,x_{n+1},\ x_n + (1+\eta_n)\,x_{n+1}\bigr)
\end{cases}
\]

Valeurs propres de $\sigma_i$ : $1$ de multiplicité $n-1$, et $\eta_i$.

Ordre de $\sigma_i \overset{?}{=}$ ordre de $\eta_i \in k^*$ si $\eta_i \ne 1$.

\section{« Étude locale » des $n$-polyèdres et des $n$-cartes (vieilles rédactions)}
\note{Titre de sa main, au crayon, sur la feuille de couverture p.~38 : « \struck{\ill{}} “étude locale” des $n$-polyèdres et des $n$-cartes (vieilles rédactions) ». Un premier mot, lourdement biffé, précède.}

\page{39}
Soit \struck{\ill{}} une \uncertain{variante} combinatoire \add{$C$} de $n$-quelette $E$ -- un $\Gamma_n$-ensemble -- où
\[
\Gamma_n = \{\sigma_0, \dots, \sigma_n \mid \sigma_0^2 = \cdots = \sigma_n^2 = 1,\ \sigma_i\sigma_j = \sigma_j\sigma_i \text{ si } j\ge i+2\}
\]
Rappelons que les $i$-facettes de la carte sont les éléments de l'ens.
\[
D_i = E/(\sigma_0, \dots, \widehat{\sigma_i}, \dots, \sigma_n) = E/\Gamma_n^{-}(i)\times\Gamma_n^{+}(i)
\]
\add{où $\Gamma_n^{-}(i) \subset \Gamma_n$, $\Gamma_n^{+}(i) \subset \Gamma_n$ sont engendrés \uncertain{resp.} par $\sigma_0 \dots \sigma_{i-1}$, \uncertain{et} $\sigma_{i+1}, \dots, \sigma_n$}.
Considérons, pour $f_i \in D_i$ fixé, l'\uncertain{ensemble} des facettes \add{$f_j$} de $C$ \struck{\uncertain{qui sont incidentes}} $< f_i$, \struck{\ill{}} i.e.\ de dim.\ $j < i$ et telles que $(f_j, f_i)$ proviennent d'un él.\ de $E$. Soit \struck{$E(i)^{+}$}
\[
E^{-}(i) = E/\Gamma_n^{-}(i) \longrightarrow D_i
\]
\add{« \uncertain{Drapeaux} \ill{} »} \struck{\ill{}}

c'est un ens.\ où opère $\Gamma_n^{+}(i) \simeq \Gamma_{i-1}$ \struck{\ill{} fibres de $D_i$ \ill{} \uncertain{compatible} \ill{} \ill{}}, et sauf erreur c'est une $(i-1)$-\uncertain{quelette}, \add{\uncertain{\ill{}}} i.e.\ $\sigma_0, \dots, \sigma_{i-1}$ et les produits \add{de} tels éléments qui commutent opèrent dans les fibres de $E^{-}(i) \to D_i$. \struck{\uncertain{Ainsi} fibre de} \add{$E^{-}(i)$} correspond : \struck{\ill{}} \struck{\ill{}} $(i-1)$-cartes combinatoires, dont les comp.\ connexes correspondent aux él.\ de $D_i$ i.e.\ aux $i$-facettes $f_i$. Soit \struck{les $(i-1)$-cartes \uncertain{combinatoires} \uncertain{quotients} \ill{}} $\Pi_{f_i}$ la fibre de $f_i$. Les $j$-facettes \uncertain{sont}
\note{Dans les énoncés de cette page, les signes $\pm$ en exposant de $\Gamma_n(i)$ et de $E(i)$ sont repassés et en partie empâtés ; ils sont lus d'après la définition « engendrés par $\sigma_0\dots\sigma_{i-1}$, et $\sigma_{i+1},\dots,\sigma_n$ » et l'isomorphisme $\Gamma_n^{+}(i)\simeq\Gamma_{i-1}$ tel qu'il est écrit, bien que ce dernier semble appeler $\Gamma_n^{-}(i)$.}

\page{40}
les él.\ de $E^{-}(i)_{f_i} / (\sigma_0, \dots, \widehat{\sigma_j}, \dots, \sigma_{i-1})$, i.e.\ les él.\ de
\[
E^{-}(i)/(\sigma_0, \dots, \widehat{\sigma_j}, \dots, \sigma_{i-1}) \simeq E/(\sigma_0, \dots, \widehat{\sigma_j}, \dots, \widehat{\sigma_i}, \dots, \sigma_n)
\]
qui \struck{\uncertain{modulo}} s'envoient dans $f_i \in D_i$, i.e.\ les drapeaux de type $(j, i)$ dont la composante de degré \uncertain{dim} $i$ est $f_i$. Plus gén., les drapeaux de type $(j_1 < \cdots < j_p)$ (où $0\le j_1 < \cdots < j_p \le i-1$) sont les drapeaux de type $(j_1, \dots, j_p, i)$ de $\Pi$ \struck{\ill{}} dont la composante de degré $i$ est $f_i$. On l'appelle \add{« \uncertain{bord} »} le « bord » de la facette $f_i$, notée $\partial f_i$.

Dualement, considérons
\[
E^{+}(i) = E/\Gamma_n^{+}(i) \longrightarrow D_i \ \overset{\text{\struck{?}}}{=}\ D_{i,i+1,\dots,n}
\]
\note{La glose « $= D_{i,i+1,\dots,n}$ », au crayon, est écrite au-dessus de la flèche.}
sur laquelle $\Gamma_n^{+}(i) \simeq \struck{\ill{}}\ \Gamma_{n-i-1}$ opère, la projection étant $D_i$. On trouve une \add{\uncertain{variante} de $(n-i-1)$-carte} $n$-\uncertain{quelette} (\uncertain{sauf} erreur), dont les facettes de dim.\ $j$ ($0\le j\le n-i-1$) « sont » \add{de $\Pi$} les drapeaux de type $(i, i+j+1)$ \struck{\uncertain{dont} la} \add{\ill{}} composante de dim.\ $i$ \struck{est un drapeau \ill{} de type $(i, j)$, \uncertain{différent} \ill{}}.

On peut aussi, pour \add{\struck{\ill{}}} \add{$j\ge i+2$} \struck{\ill{}}
\[
f_i < f_j
\]
considérer \struck{$E(i,j) = E/(\sigma_0 \dots \sigma_{i-1};\ \sigma_{j+1} \dots \sigma_n)$}
\note{La page s'arrête sur « considérer » ; la suite, si elle existe, est au-delà du lot. Le bas de la page est très chargé d'interlignes et de ratures, et la syntaxe exacte des deux dernières phrases reste incertaine.}

\end{document}
